\documentclass[options]{philosophersimprint}
%\usepackage{sectsty}
%\linespread{1.2}
\usepackage{xunicode}
\usepackage{xltxtra}
\defaultfontfeatures{Mapping=tex-text}
\setmainfont[
Ligatures={Common}, Numbers={OldStyle}, Variant=01,
BoldFont=LinLibertine_RB.otf,
ItalicFont=LinLibertine_RI.otf,
BoldItalicFont=LinLibertine_RBI.otf
]{LinLibertine_R.otf}
\setmonofont[Scale=0.8]{DejaVuSansMono.ttf}
\usepackage{stmaryrd}
\usepackage[toc,page]{appendix}
\usepackage{amsmath,amsthm,amssymb}
\usepackage[colorlinks=true, citecolor=blue, linkcolor=blue, hyperfootnotes=false]{hyperref}
\usepackage{nameref}
\usepackage{soul}
\usepackage{frame}
\usepackage[style=apa, natbib=true]{biblatex}
\usepackage{scrextend}
\usepackage{mathrsfs}
%\addtokomafont{labelinglabel}{\sffamily}
\makeatletter
\newcommand\labelitem[2]{%
  \item[#1]\def\@currentlabel{#1}\label{#2}}
\makeatother
\usepackage[inline]{enumitem}
\usepackage{framed, csquotes}
\usepackage{tikz}
\usepackage{scalerel}
\usepackage{titlesec}
\newcommand\halo{{{\circ}\mkern-1mu}}
\usepackage{linguex}
\usepackage{parskip}
\usepackage{fancyhdr}
\pagestyle{plain}
% Diagrams
%\usepackage{pgfplots}
\usepackage{tikz}
\usepackage{qtree}
\usepackage{comment}



\usepackage{FOL_Yale/notation}

\begin{document}

PHIL 1115: First-order Logic \hfill Yale, Fall 2026\\
Lecture 7, Handout \hfill Dr. Daniel Grimmer

\textbf{1. Invalidity Via Trees} In the previous lecture, I gave definitions for tree-contradiction, tree-tautology, tree-contingency, and tree-equivalence. But what about tree-validity and tree-invalidity? We turn to these now. Your friend claims that a certain argument (with premises, $X$, and conclusion, $A$) is valid, but you do not believe them. How can you use a tree to show that their argument is invalid?

Recall (from Lecture 4) how we could invalidate arguments using the table method: we constructed a (potentially huge) truth table and looked for a \textit{countermodel}. That is, we looked for a model (a row of the truth table) in which the argument's set of premises is true and yet its conclusion is false. Now recall what trees do: \textbf{they search for ways of simultaneously satisfying a whole list of formulas.} If we can find a model according to which $X$ and $\Neg A$ are all true, this will produce a counterexample to their argument. We are thus led to the following definition.

\textbf{Invalidity:} An argument with premises $X$ and conclusion $A$ is \textbf{tree-invalid} (denoted $X \not\Proves A$) if and only if its premises together with the negation of its conclusion are jointly tree-satisfiable (denoted $X,\Neg A\notProves$). To test if an argument is tree-invalid, one develops the tree for $X$ and $\Neg A$ to check if it has an open branch. Importantly, this is a different notion from table-invalidity which we denoted $X \notEntails A$ in Lecture 4.

Exercise: Consider the argument with premises $X: \  \Neg(p\Disj q) \Disj (r \Conj \Neg s)$ and $s \Cond (r \Disj q)$ and with conclusion $A: \ p \Cond s$. Use a tree to show that this argument is tree-invalid.

The fact that the following tree remains open means that there is a model (i.e., a logical possibility) in which its roots are all true. In this world, the argument's premises, $X$, are all true but its conclusion, $A$, is false (because $\Neg A$ is true). The argument has a countermodel, and is hence invalid. We can denote this conclusion as
$$\Neg(p\Disj q) \Disj (r \Conj \Neg s), \quad s \Cond (r \Disj q)\quad \not\Proves \quad p \Cond s$$
Had we demonstrated this result using the table method, we would instead write, $\Neg(p\Disj q) \Disj (r \Conj \Neg s), \quad s \Cond (r \Disj q)\quad \notEntails \quad p \Cond s$.

\begin{center}
\Tree[.{
\text{Start from $X$ and $\Neg A$:} \\
.\\
$X:\quad \Neg(p\Disj q)\Disj(r\Conj\Neg s)$\quad\checkmark \\ \qquad$s\Cond(r\Disj q)$\quad\checkmark \\ $\Neg A:\quad \Neg(p\Cond s)$\quad\checkmark \\ $\vert$ \\ $p$ \\ $\Neg s$} [.{\qquad$\Neg(p\Disj q)$\quad\checkmark \\ $\vert$ \\ $\Neg p$ \\ $\Neg q$ \\ x} ] [.{\qquad$r\Conj\Neg s$\quad\checkmark \\ $\vert$ \\ $r$ \\ $\Neg s$} [.{$\Neg s$ \\ o} ] [.{\qquad$r\Disj q$\quad\checkmark} [.{$r$ \\ o} ] [.{$q$ \\ o} ] ] ] ]
\end{center}

Importantly, the tree-test for invalidity is constructive; each open branch highlights at least one model (i.e., a logical possibility) in which the argument fails. \textbf{We shall prove this claim later on in today's lecture.} In this example we can read a model off of the atomic (and negated-atomic) formulas on the right-most branch. From top to bottom we have: $p$ and $\Neg s$ and $r$ and $\Neg s$ and $q$. This open branch specifies a countermodel to the argument:
\begin{align}\label{EqEval}
I: \quad p=T, \quad q=T, \quad  r=T, \quad  s=F
\end{align}
If you were to build a 16-row truth table for this argument (see Lecture 4), and look at this row, then you would find that the argument's premises are all true and yet its conclusion is false. More on this later.

\textbf{2. Overview: Soundness and Completeness for Trees}

In the next two lectures we will prove the \textbf{soundness} and \textbf{completeness} of the tree method.

\textbf{Soundness}. The tree method is \textbf{sound}: if $X \Proves A$, then $X \Entails A$ (i.e., the tree method recognizes as valid \textit{only} arguments that the truth table method also considers valid).

\textbf{Completeness}. The tree method is \textbf{complete}: if $X \Entails A$, then $X \Proves A$ (i.e., the tree method recognizes as valid \textit{all} arguments that the truth table method also recognizes as valid).

Together, these results tell us that $X \Proves A$ if and only if $X \Entails A$. That is, if you had one computer that could make truth tables and another one that could make truth trees, you could trust that they would always reach the exact same conclusions about the validity of arguments.

None of this is trivial. Lots of ways of altering our system of tree rules would ruin soundness, completeness, or both. \textbf{Ask your TFs about this!}

In this lecture we shall (sketch) a proof of the completeness of the tree method. Before doing this, however, we should first rephrase it using the \textit{contrapositive rule}. In Lecture 4, we used the table method to prove that $p\Cond q$ is logically equivalent to $\Neg q\Cond \Neg p$. Moreover, in Lecture 6 we proved this using the tree method. Elevating this rule to our meta-language (English) and applying it to the above formulation of completeness (if $X \Entails A$, then $X \Proves A$) gives us the following alternative formulation.

\textbf{Proposition 1} (Completeness of the Tree Method)\textbf{.} Let $X$ be a set of formulas and let $A$ be a formula. Consider the argument with premises given by $X$ and conclusion given by $A$. If this argument is tree-invalid, $X \not\Proves A$, then it is also table-invalid, $X \not\Entails A$.

Before proving this claim, let us work through an example in detail.

\textbf{3. A Worked Example: From Open Branch to Satisfying Model}

Recall from our invalid example how we could read a model, $I$, off of any open branch of our tree. Concretely, recall that looking at its right-most branch (the one ending in $q$) gave us the following atomic and negated-atomic formulas:
$$p \text{ and } \Neg s \text{ and } r \text{ and } \Neg s \text{ and } q.$$
Note that these formulas are all jointly satisfiable. Namely the model,
$$I: \quad p=T, \quad q=T, \quad  r=T, \quad  s=F$$
which we found in Eq.\eqref{EqEval} makes all of these formulas true. This is not a coincidence. Here is a general fact.

\textbf{Lemma 1} (Open Branch Implies Atomic Satisfaction)\textbf{.} Let $O$ be an open branch on some completely resolved tree. Consider all of the atomic formulas on $O$ as well as all of the negations of atomic formulas on $O$. \textit{Claim}: There exists some model, $I$, according to which all of these atomic and negated-atomic formulas are true.

\begin{proof}
For any such open branch we can construct a model, $I$, as follows. To each of the atomic formulas (e.g., $p$) which appear on $O$, our model $I$ shall assign a truth value of $T$ (such that $p$ is true). Similarly, to each of the negated atomic formulas (e.g., $\Neg s$) which appear on $O$, our model $I$ shall assign the underlying atomic formula a truth value of $F$ (e.g., $s$ false, so that $\Neg s$ is true). Any atomic formulas not appearing on $O$ may be assigned arbitrarily. These two prescriptions do not conflict since (by assumption) the branch $O$ is open, and hence it doesn't contain any contradictory atomic formulas. By construction, all of these atomic and negated-atomic formulas are true according to this model.
\end{proof}

But here is the key idea we're after: it's not just the \textit{atomic} (and negated-atomic) formulas on this branch that are true according to $I$; \textit{every} formula on this branch is true according to $I$, all the way up to the roots. Let's check this, working from the bottom of the branch to the top, one resolution step at a time:
\begin{itemize}
    \item $q$ is true according to $I$: this is just one of our atomic assignments.
    \item $C_1 \ := \ r\Disj q$ is also true according to $I$. Why? We can check it directly as follows: a disjunction is true whenever at least one of its disjuncts is true. Note that $q$ is true according to $I$ and hence so is $r\Disj q$.
    \item Moving up the open branch, we find $\Neg s$ and then $r$. These are both true according to $I$ by construction (we used these formulas to build $I$).
    \item Next we find $C_2 \ := \ r\Conj\Neg s$. A conjunction is true if and only if both conjuncts are true. Since we just established that $r$ and $\Neg s$ are both true according to $I$, so too is $r\Conj\Neg s$.
    \item Moving up the open branch, we find $\Neg s$ and then $p$. These are both true according to $I$ by construction (we used these formulas to build $I$).
    \item Next we find $C_3 \ := \ \Neg(p\Cond s)$. A negated conditional is true if and only if its antecedent is true but its consequent is false. Since we just established that $p$ and $\Neg s$ are both true according to $I$, so too is $\Neg(p\Cond s)$.
    \item At last we get to a more difficult one. Why is $C_4 \ := \ s\Cond (r \Disj q)$ true according to $I$? Unlike all of the previous complex formulas, the effect that resolving $C_4$ had on the tree is not immediately downstream. When we resolved $C_4$ it caused the split later on between $\Neg s$ and $r\Disj q$. Because $C_4$ is on the open branch, $O$, we know that one half of this split must also be on $O$ (namely, the $r\Disj q$ branch is on $O$). Because $r\Disj q$ sits lower upon $O$ than $C_4$ does, it must have been covered already by our above discussion. Hence, it must be true according to $I$. But this is sufficient to guarantee that $C_4$ itself is true according to $I$. (For all splitting rules, the parent is true according to $I$ so long as at least one of its two branches is).
    \item Lastly, we have $C_5 \ := \ \Neg(p\Disj q)\Disj(r\Conj\Neg s)$. Why is this formula true according to $I$? The reason is almost exactly the same as for $C_4$. When we resolved $C_5$ it caused the split later on between $\Neg(p\Disj q)$ and $r\Conj\Neg s$. Because $C_5$ is on the open branch, $O$, we know that one half of this split must also be on $O$ (namely, the $r\Conj\Neg s$ branch is on $O$). Because $r\Conj\Neg s$ sits lower upon $O$ than $C_5$ does, it must have been covered already by our above discussion. Hence, it must be true according to $I$. But this is sufficient to guarantee that $C_5$ itself is true according to $I$. (Again, for all splitting rules, the parent is true according to $I$ so long as at least one of its two branches is).
\end{itemize}

So every formula, \textit{and not just the (negated) atoms}, on this open branch, $O$, comes out true according to $I$. Importantly, this includes the roots of the tree: namely, the argument's premises $X$ and its negated conclusion, $\Neg A$. This explicitly demonstrates that $I$ is a genuine countermodel to the argument: it makes both premises of our argument true while making the conclusion false. We've just shown, very concretely, that $X \not\Proves A$ yields $X \notEntails A$ in this particular case.

\textbf{4. Proof of Completeness (Sketch)} We just walked through the completeness of the tree method by hand, for one branch of one tree. Let's now see why it always holds, for any complete tree whatsoever. Hopefully it is clear that the techniques we used in this particular example generalize. Notice the pattern of reasoning:
\begin{enumerate}
    \item $X \not\Proves A$ gives us an open branch, $O$.
    \item From this open branch, Lemma 1 then gives us a model $I$ according to which all of the (negated) atomic formulas on $O$ are true.
    \item Some more argumentation (generalized as Lemma 2 below) then shows that \textit{every formula} on $O$ is true according to $I$, even the complex ones.
    \item In particular, the roots of the tree are all true according to $I$. Because we chose the roots to be $X$ and $\Neg A$, this constitutes a countermodel to the argument. Hence we have $X \notEntails A$.
\end{enumerate}

Once we generalize the third step in this reasoning, we will have shown \textit{in general} that $X \not\Proves A$ implies $X \notEntails A$. Namely, whenever the tree method deems an argument invalid, so does the table method. This is our completeness result.

The following Lemma generalizes step 3 using a proof by induction.

\textbf{Lemma 2} (Atomic Satisfaction Implies Branch Satisfaction)\textbf{.} Let $O$ be some branch on some completely resolved tree with the following property. There exists some model, $I$, according to which all of the atomic and negated-atomic formulas on $O$ are true. Then every formula on $O$ is true according to $I$.
\begin{proof}
Consider the sequence of complex formulas, $C_n$, numbered from bottom to top up the given branch, $O$, of this tree. We shall proceed with a proof by induction on this index $n$. The property to be established at each step in this construction is ``being true according to $I$''. We shall prove the inductive step before the base case. (Either order works: neither half uses the other.)

\textit{Inductive Step}. Consider any complex formula, $C_{n+1}$, on branch $O$. Our inductive hypothesis is that all of the complex formulas on $O$ below this point (i.e., from $C_n$ downwards) are true according to $I$. Our goal is to use this assumption to show that $C_{n+1}$ itself is true according to $I$.

$C_{n+1}$ must be of one of the nine forms for which we have defined tree rules. Because our tree is complete, $C_{n+1}$ has been resolved by applying one of these nine resolution rules. Whenever $C_{n+1}$ was resolved it introduced some new descendant formulas lower down on $O$ (some tree rules have one descendant per branch, and others have two per branch). Whatever these descendants on $O$ are, they must be either atomic formulas, negated-atomic formulas, or complex formulas lower on $O$. However, in any case, the descendants of $C_{n+1}$ on $O$ will be true according to $I$. For all of the atomic and negated-atomic formulas on $O$ this is true by construction (indeed, we built $I$ out of them). For complex descendants of $C_{n+1}$ on $O$, we know they are true according to $I$ because of our inductive hypothesis.

In sum: for every tree-rule, all of the descendants of $C_{n+1}$ on $O$ are all true according to $I$. What we shall show next is that in all nine cases this is sufficient to establish that $C_{n+1}$ itself is true according to $I$.

\textit{Conjunction}. Suppose that $C_{n+1}= A \Conj B$ is a conjunction. By resolving this conjunction, $A$ and $B$ will both appear somewhere on $O$. Hence, by our above reasoning, both $A$ and $B$ will be true according to $I$. From this it then follows that $C_{n+1}= A \Conj B$ is also true according to $I$.

\textit{Negated Conjunction}. Suppose that $C_{n+1}= \ \Neg(A \Conj B)$ is a negated conjunction. Resolving this negated conjunction will result in a split forming along $O$ somewhere below $C_{n+1}$. One of these branches will contain $\Neg A$ and the other shall contain $\Neg B$. Either of these two continuations could be the $O$ branch which we are interested in.

If $O$ is the branch containing $\Neg A$ then our above reasoning tells us that $\Neg A$ is true according to $I$. From this it would then follow that $C_{n+1}= \Neg(A \Conj B)$ is also true according to $I$.

If $O$ is the branch containing $\Neg B$ then our above reasoning tells us that $\Neg B$ is true according to $I$. From this it would then follow that $C_{n+1}= \Neg(A \Conj B)$ is also true according to $I$.

The other seven cases follow the same structure.

\textit{Base Case}. For our base case, consider $C_1$, the lowest complex formula on branch $O$. Much of the above reasoning applies in this case as well. The key difference is that because $C_1$ is the \textit{lowest complex} formula on $O$, its descendants on $O$ must be atomic or negated-atomic (they cannot be complex). But by construction we know that all such formulas on $O$ are true according to $I$. Repeating the case-by-case analysis from above then establishes that $C_1$ is true according to $I$.
\end{proof}
One can combine Lemmas 1 and 2 to prove Proposition 1 (Completeness of the Tree Method) as we sketched out above.

\end{document}